\section*{Capitulo 2. Polinomios y fracciones algebraicas}

\begin{shortanswer}{[GX-C02.S01-01]}
\[
Q(3)=-2.
\]
\end{shortanswer}

\begin{shortanswer}{[GX-C02.S01-02]}
\[
R(-1)=0.
\]
\end{shortanswer}

\begin{shortanswer}{C02.S01 practica autonoma}
\begin{enumerate}[label=\textbf{PX-C02.S01-\arabic*:}, leftmargin=*]
  \item \(5\)
  \item \(5\)
  \item \(0\)
  \item \(25\)
  \item \(2\)
  \item \(5\)
\end{enumerate}
\end{shortanswer}

\begin{shortanswer}{[GX-C02.S02-01]}
\[
(x+5)^2=x^2+10x+25.
\]
\end{shortanswer}

\begin{shortanswer}{[GX-C02.S02-02]}
\[
(3x-2)(3x+2)=9x^2-4.
\]
\end{shortanswer}

\begin{shortanswer}{C02.S02 practica autonoma}
\begin{enumerate}[label=\textbf{PX-C02.S02-\arabic*:}, leftmargin=*]
  \item \(x^2-8x+16\)
  \item \(4x^2+4x+1\)
  \item \(x^2-9\)
  \item \(x^3+x^2+x+6\)
  \item \(2x^2-9x+4\)
  \item \(4x\)
\end{enumerate}
\end{shortanswer}

\begin{shortanswer}{[GX-C02.S03-01]}
Cociente: \(x^2-5x+6\). Resto: \(0\).
\end{shortanswer}

\begin{shortanswer}{[GX-C02.S03-02]}
Cociente: \(3x^2-5x\). Resto: \(8\).
\end{shortanswer}

\begin{shortanswer}{C02.S03 practica autonoma}
\begin{enumerate}[label=\textbf{PX-C02.S03-\arabic*:}, leftmargin=*]
  \item Cociente: \(x-3\). Resto: \(0\).
  \item Cociente: \(2x-1\). Resto: \(-2\).
  \item Cociente: \(x^2+4x+4\). Resto: \(0\).
  \item Cociente: \(4x^2+4x+4\). Resto: \(0\).
  \item Cociente: \(x^2-x+1\). Resto: \(-2\).
  \item Cociente: \(2x^2-x+3\). Resto: \(-12\).
\end{enumerate}
\end{shortanswer}

\begin{shortanswer}{[GX-C02.S04-01]}
Resto: \(4\).
\end{shortanswer}

\begin{shortanswer}{[GX-C02.S04-02]}
\[
b=-5.
\]
\end{shortanswer}

\begin{shortanswer}{C02.S04 practica autonoma}
\begin{enumerate}[label=\textbf{PX-C02.S04-\arabic*:}, leftmargin=*]
  \item \(9\)
  \item Si
  \item \(a=-2\)
  \item \(0\)
  \item \(c=3\)
  \item Si
\end{enumerate}
\end{shortanswer}

\begin{shortanswer}{[GX-C02.S05-01]}
\[
x^2-9=(x-3)(x+3).
\]
\end{shortanswer}

\begin{shortanswer}{[GX-C02.S05-02]}
\[
x^4-5x^2+4=(x-1)(x+1)(x-2)(x+2).
\]
\end{shortanswer}

\begin{shortanswer}{C02.S05 practica autonoma}
\begin{enumerate}[label=\textbf{PX-C02.S05-\arabic*:}, leftmargin=*]
  \item \((x+2)(x+3)\)
  \item \(x(x-2)^2\)
  \item \((x-1)(x+1)\)
  \item \(2(x-2)(x+2)\)
  \item \(x^2(x+3)\)
  \item \((x-2)(x+2)(x^2+4)\)
\end{enumerate}
\end{shortanswer}

\begin{shortanswer}{[GX-C02.S06-01]}
\[
\frac{x^2-4}{x^2-2x}=\frac{x+2}{x},\qquad x\neq 0,2.
\]
\end{shortanswer}

\begin{shortanswer}{[GX-C02.S06-02]}
\[
\frac{x^2-1}{x^2+2x+1}=\frac{x-1}{x+1},\qquad x\neq -1.
\]
\end{shortanswer}

\begin{shortanswer}{C02.S06 practica autonoma}
\begin{enumerate}[label=\textbf{PX-C02.S06-\arabic*:}, leftmargin=*]
  \item \(D=\mathbb{R}\setminus\{4\}\)
  \item \(\dfrac{x}{x+1}\), con \(x\neq -1,1\)
  \item \(\dfrac{x+4}{x-4}\), con \(x\neq 4\)
  \item \(x+5\), con \(x\neq 0\)
  \item \(-(x+5)\), con \(x\neq 5\)
  \item \(\dfrac{x-3}{x+3}\), con \(x\neq -3,3\)
\end{enumerate}
\end{shortanswer}

\begin{shortanswer}{[GX-C02.S07-01]}
\[
\frac{1}{x}+\frac{1}{2x}=\frac{3}{2x}.
\]
\end{shortanswer}

\begin{shortanswer}{[GX-C02.S07-02]}
\[
\frac{2}{x+2}.
\]
\end{shortanswer}

\begin{shortanswer}{C02.S07 practica autonoma}
\begin{enumerate}[label=\textbf{PX-C02.S07-\arabic*:}, leftmargin=*]
  \item \(\dfrac{3}{x}\)
  \item \(\dfrac{2x}{x^2-4}\)
  \item \(\dfrac{5}{2x}\)
  \item \(\dfrac{x-3}{x^2-1}\)
  \item \(1\)
  \item \(-\dfrac{x}{x^2-1}\)
\end{enumerate}
\end{shortanswer}

\begin{shortanswer}{[GX-C02.S08-01]}
\[
\frac{x}{2}.
\]
\end{shortanswer}

\begin{shortanswer}{[GX-C02.S08-02]}
\[
\frac{x-2}{x}.
\]
\end{shortanswer}

\begin{shortanswer}{C02.S08 practica autonoma}
\begin{enumerate}[label=\textbf{PX-C02.S08-\arabic*:}, leftmargin=*]
  \item \(\dfrac{x-1}{x+1}\)
  \item \(\dfrac{x+3}{x}\)
  \item \(\dfrac{x+2}{2}\)
  \item \(1\)
  \item \(\dfrac{x+1}{x}\)
  \item \(\dfrac{2(x-1)}{3}\)
\end{enumerate}
\end{shortanswer}

\begin{shortanswer}{[GX-C02.S09-01]}
\[
x=2.
\]
\end{shortanswer}

\begin{shortanswer}{[GX-C02.S09-02]}
\[
x=\frac{7}{2}.
\]
\end{shortanswer}

\begin{shortanswer}{C02.S09 practica autonoma}
\begin{enumerate}[label=\textbf{PX-C02.S09-\arabic*:}, leftmargin=*]
  \item \(x=\dfrac{1}{5}\)
  \item \(x=2\)
  \item \(x=2\)
  \item \(x=1\)
  \item Sin solucion
  \item \(x=2\)
\end{enumerate}
\end{shortanswer}
